% Parkinson's Disease Voice Classification - MSc Thesis
% Preamble: All packages, configuration, and custom commands

%% ============================================================================
%% DOCUMENT METADATA
%% ============================================================================

% Author information
\newcommand{\myAuthor}{Μάριος Χρήστος Συρμακέζης}
\newcommand{\myTitle}{Voice-Based Classification of Parkinson's Disease Using Classical Machine Learning}
\newcommand{\myTitleLineBreak}{Voice-Based Classification of Parkinson's Disease Using Classical Machine Learning}
\newcommand{\myTitleGR}{Ταξινόμηση Νόσου Parkinson μέσω Φωνής με Χρήση Κλασικής Μηχανικής Μάθησης}

% Supervisor information
\newcommand{\mySupervisor}{Γεώργιος Ματσόπουλος}
\newcommand{\mySupervisorTitle}{Καθηγητής ΗΜΜΥ Ε.Μ.Π.}

% Committee members
\newcommand{\myExaminerA}{Ουρανία Πετροπούλου}
\newcommand{\myExaminerATitle}{Επίκουρη Καθηγήτρια ΗΜΜΥ Ε.Μ.Π.}
\newcommand{\myExaminerB}{[Examiner 2 Name]}
\newcommand{\myExaminerBTitle}{[Examiner 2 Title]}

% Date and location
\newcommand{\myDate}{Φεβρουάριος 2026}
\newcommand{\myLocation}{Αθήνα}
\newcommand{\myApprovalDate}{[26 Φεβρουαρίου 2026]}

% Institution
\newcommand{\myUniversity}{Εθνικό Μετσόβιο Πολυτεχνείο}
\newcommand{\mySchool}{Σχολή Ηλεκτρολόγων Μηχανικών και Μηχανικών Η/Υ}
\newcommand{\myDivision}{Συστημάτων Μετάδοσης Πληροφορίας \& Τεχνολογίας Υλικών}

%% ============================================================================
%% PACKAGES
%% ============================================================================

% Encoding and fonts (XeLaTeX)
\usepackage{fontspec}
\setmainfont{FreeSerif}
\setsansfont{FreeSans}
\setmonofont{FreeMono}

% Language
\usepackage{polyglossia}
\setdefaultlanguage{greek}
\setotherlanguage{english}

% Greek captions for structural elements
% NOTE: The class file (ntua-thesis.cls) loads babel with [greek,english],
% making English the active language. We must override BOTH caption sets
% to ensure Greek labels appear regardless of which language babel considers active.
\addto\captionsgreek{%
  \renewcommand{\contentsname}{Πίνακας Περιεχομένων}%
  \renewcommand{\listfigurename}{Κατάλογος Σχημάτων}%
  \renewcommand{\listtablename}{Κατάλογος Πινάκων}%
  \renewcommand{\chaptername}{Κεφάλαιο}%
  \renewcommand{\bibname}{Βιβλιογραφία}%
  \renewcommand{\appendixname}{Παράρτημα}%
}
\addto\captionsenglish{%
  \renewcommand{\contentsname}{Πίνακας Περιεχομένων}%
  \renewcommand{\listfigurename}{Κατάλογος Σχημάτων}%
  \renewcommand{\listtablename}{Κατάλογος Πινάκων}%
  \renewcommand{\chaptername}{Κεφάλαιο}%
  \renewcommand{\bibname}{Βιβλιογραφία}%
  \renewcommand{\appendixname}{Παράρτημα}%
}

% Page layout (matches NTUA template)
\usepackage[
    top=3cm,
    bottom=3cm,
    left=2.5cm,
    right=2.5cm,
    bindingoffset=0.5cm,
    headheight=15pt
]{geometry}

% Headers and footers
\usepackage{fancyhdr}

% Graphics and figures
\usepackage{graphicx}
\usepackage{float}
\usepackage{subcaption}

% Tables
\usepackage{booktabs}
\usepackage{longtable}
\usepackage{multirow}
\usepackage{array}

% Math
\usepackage{amsmath}
\usepackage{amssymb}
\usepackage{amsthm}

% Algorithms
\usepackage{algorithm}
\usepackage{algpseudocode}

% Code listings
\usepackage{listings}
\usepackage{xcolor}

% Hyperlinks
\usepackage[
    colorlinks=true,
    linkcolor=blue,
    citecolor=blue,
    urlcolor=blue,
    bookmarksdepth=3,
]{hyperref}
\usepackage{bookmark}

% Bibliography
\usepackage[numbers,sort&compress]{natbib}

% Miscellaneous
\usepackage{enumitem}
\usepackage{parskip}
\usepackage{setspace}

% Chapter heading style
\usepackage{titlesec}

%% ============================================================================
%% CUSTOM COMMANDS
%% ============================================================================

% Blank page utility (used after front-matter chapters, per NTUA template)
\def\blankpage{%
      \clearpage%
      \thispagestyle{empty}%
      \null%
      \clearpage}

\raggedbottom

%% ============================================================================
%% CONFIGURATION
%% ============================================================================

% Line spacing
\onehalfspacing

% Code listing colors
\definecolor{codegreen}{rgb}{0,0.6,0}
\definecolor{codegray}{rgb}{0.5,0.5,0.5}
\definecolor{codepurple}{rgb}{0.58,0,0.82}
\definecolor{backcolour}{rgb}{0.95,0.95,0.92}

% Code listing style
\lstdefinestyle{mystyle}{
    backgroundcolor=\color{backcolour},
    commentstyle=\color{codegreen},
    keywordstyle=\color{magenta},
    numberstyle=\tiny\color{codegray},
    stringstyle=\color{codepurple},
    basicstyle=\ttfamily\footnotesize,
    breakatwhitespace=false,
    breaklines=true,
    captionpos=b,
    keepspaces=true,
    numbers=left,
    numbersep=5pt,
    showspaces=false,
    showstringspaces=false,
    showtabs=false,
    tabsize=2,
    frame=single
}
\lstset{style=mystyle}

% Python language style for code listings
\lstdefinelanguage{Python}{
    keywords={def, class, return, if, elif, else, for, while, import, from, as, with, try, except, raise, True, False, None, and, or, not, in, is, lambda, yield, pass, break, continue, global, nonlocal, assert, del},
    keywordstyle=\color{blue}\bfseries,
    ndkeywords={self, __init__, __name__, print, range, len, int, str, float, list, dict, set, tuple},
    ndkeywordstyle=\color{magenta},
    sensitive=true,
    comment=[l]{\#},
    morecomment=[s]{'''}{'''},
    morecomment=[s]{"""}{"""},
    commentstyle=\color{codegreen}\itshape,
    stringstyle=\color{codepurple},
    morestring=[b]',
    morestring=[b]"
}

% Header style
\pagestyle{fancy}
\fancyhf{}
\fancyhead[LE,RO]{\thepage}
\fancyhead[RE]{\nouppercase{\leftmark}}
\fancyhead[LO]{\nouppercase{\rightmark}}
\renewcommand{\headrulewidth}{0.4pt}

% Chapter heading style
\titleformat{\chapter}[display]
    {\normalfont\huge\bfseries}
    {\chaptertitlename\ \thechapter}
    {20pt}
    {\Huge}

% Graphics path
\graphicspath{{figures/}{../outputs/plots/}{../assets/images/}}
